% Reproduced inherited proof; finding P-F19.
\subsection{An exponential canonical modular index inside ordinary plethysm}

\begin{proposition}[P-F19: chain cubics obstruct the canonical clause of PH3]
\label{prop:cubiccanonicalcounterexample}
Let $n\ge2$, put
\[
 g=\gcd(n,3),\qquad d_0=\frac ng,\qquad
 W=\operatorname{Sym}^3\mathbb C^n,\qquad
 m=\dim W=\binom{n+2}{3},
\]
and consider the ordinary plethysm ray with
\[
 \mu=(3),\qquad \lambda=(d_0),\qquad
 \pi=((3/g)^n).
 \tag{P-F19a}
\]
Its canonical multiplicity algebra is
\[
 T^{(n)}=\bigoplus_{s\ge0}
  \operatorname{Hom}_{GL_n}
  \left((\det)^{3s/g},\operatorname{Sym}^{d_0s}W\right),
 \tag{P-F19b}
\]
with the Cartan multiplication. Every homogeneous system of parameters
of $T^{(n)}$ contains a parameter of stretch degree at least
$2^{n-1}$. In particular its canonical modular index satisfies
\[
 \chi(T^{(n)})\ge2^{n-1}.
 \tag{P-F19c}
\]
The parameter of GCT VI, Hypothesis 1.6.7, is
\[
 h=m^2+\operatorname{ht}(\lambda)+\operatorname{ht}(\pi)
   =\binom{n+2}{3}^{\!2}+n+1=O(n^6).
 \tag{P-F19d}
\]
Consequently the canonical modular-index clause of that hypothesis
is false already for ordinary plethysm with the fixed inner shape
$\mu=(3)$. This statement concerns the h.s.o.p. of the actual canonical algebra.
The parameter argument below alone does not assert that every positive
rational expression for its Hilbert series has a large denominator
weight. P-F31 later proves that stronger conclusion using the full
support theorem and an actual Hilbert-series pole. The superpolynomial positivity and saturation indices
require the paired-wave argument and quantitative estimates below.
\end{proposition}

\begin{proof}
Write $V=\mathbb C^n$ and identify
\[
 R=\mathbb C[W^*]^{SL_n}=(\operatorname{Sym}W)^{SL_n}.
 \tag{P-F19f}
\]
Thus points of $W^*=\operatorname{Sym}^3V^*$ are cubic forms in
the linear coordinate functions $x_1,\ldots,x_n$ on $V$.
This choice avoids identifying a representation with its dual.
The usual action on such a form is $(a\cdot F)(x)=F(a^{-1}x)$.
The ring $R$ has its ordinary coefficient grading, denoted $R_D$.

First we check the canonical grading and multiplication in (P-F19b).
If $R_D\ne0$, a scalar $\zeta I_n\in SL_n$ with $\zeta^n=1$
acts on $\operatorname{Sym}^D W$ by $\zeta^{3D}$, whence
$n\mid3D$. This is equivalent to $D=d_0s$ for some integer
$s\ge0$. On $R_{d_0s}$ the entire $GL_n$ acts by the character
$(\det)^{3s/g}$. Indeed, write $a=zh$ with $h\in SL_n$ and
$z^n=\det a$. Since $h$ fixes this subspace and $zI_n$ acts by
$z^{3d_0s}$, its action is $(\det a)^{3s/g}$, independently of
the choice of the $n$th root $z$.

Choose a nonzero volume vector $\varepsilon\in\bigwedge^nV$.
Evaluation of an intertwiner at $\varepsilon^{\otimes3s/g}$
identifies the degree-$s$ space in (P-F19b) with $R_{d_0s}$.
The Cartan product of these one-dimensional determinant modules
sends the compatible volume vectors to their tensor product.
Evaluation consequently respects multiplication in
$\operatorname{Sym}W$. We obtain an isomorphism of graded
algebras
\[
 T^{(n)}\simeq\bigoplus_{s\ge0}R_{d_0s},
 \qquad \deg_{T^{(n)}} R_{d_0s}=s.
 \tag{P-F19g}
\]
All nonzero ordinary degrees of $R$ occur in this display, so this
is the whole invariant ring with its grading divided by $d_0$,
not a proper Veronese subring. The canonical ring and its
multiplication are those in GCT VI, Section 4.1.2, equation (4.5),
printed p.~69, and equations (4.2)--(4.3), printed p.~67.
For the canonical modular index we use the detailed $T$-ring
definition of Section 4.1.1, printed p.~68.

Consider the explicit point
\[
 F_n=x_1^2x_2+x_2^2x_3+\cdots+x_{n-1}^2x_n+x_n^3
       \in W^*.
 \tag{P-F19h}
\]
Its projective hypersurface is smooth. To see this directly, its
partial derivatives are
\[
 \partial_{x_1}F_n=2x_1x_2,\qquad
 \partial_{x_i}F_n=x_{i-1}^2+2x_ix_{i+1}\ (1<i<n),\qquad
 \partial_{x_n}F_n=x_{n-1}^2+3x_n^2.
\]
Suppose they all vanish at a nonzero vector, and let $j$ be the
smallest index with $x_j\ne0$. If $j<n$, the $j$th derivative
gives $x_{j+1}=0$, and the $(j+1)$st derivative is then $x_j^2$,
which is nonzero, a contradiction. If $j=n$, the last derivative
is $3x_n^2\ne0$, also a contradiction. Thus the partials have
no common nonzero zero, proving smoothness over $\mathbb C$.

Let $\Delta$ be the discriminant of cubic forms. It is an
$SL_n$-invariant homogeneous polynomial of coefficient degree
$n2^{n-1}$, and it is nonzero precisely on smooth projective
hypersurfaces. An exact primary locator for these properties is
McCrory--Shifrin--Varley, \emph{Siegel modular forms generated by
invariants of cubic hypersurfaces}, arXiv:alg-geom/9403005v3,
Section 3.1.2(i), printed p.~15: the discriminant is the resultant
of the partial derivatives, with degree $n(k-1)^{n-1}$ for forms
of degree $k$.
For a direct resultant treatment see also Bus\'e--Jouanolou,
\emph{A computational approach to the discriminant of homogeneous
polynomials}, arXiv:1210.4697v1, Definition 4.6 and
Proposition 4.7(i), printed p.~40, and Proposition 4.13,
printed pp.~42--43, for the coordinate-change character.
Hence
\[
 \Delta(F_n)\ne0,\qquad
 \Delta\in R_{n2^{n-1}}
      =T^{(n)}_{g2^{n-1}}.
 \tag{P-F19i}
\]
In particular $F_n$ is outside the null cone. No claim about
stability beyond this nonvanishing is needed.

We next prove the decisive divisibility statement:
\[
 P\in T^{(n)}_s,\quad P(F_n)\ne0
       \quad\Longrightarrow\quad 2^{n-1}\mid s.
 \tag{P-F19j}
\]
We give two independent calculations.

\emph{Coefficient-weight calculation.}
Write $D=d_0s$ for the ordinary coefficient degree. A polynomial
invariant is a sum of coefficient monomials of torus weight zero:
the diagonal torus acts diagonally on the monomial basis, so a
nonzero weight cannot occur in an invariant polynomial.
If $P(F_n)\ne0$, at least one monomial occurring in $P$ uses
only the $n$ coefficients displayed in (P-F19h), because every
other coefficient is zero at $F_n$. Let its nonnegative exponent
tuple be $(a_1,\ldots,a_n)$; then $\sum_i a_i=D$.
Triviality of its character on the determinant-one diagonal torus
means its exponent vector in $x_1,\ldots,x_n$ is constant.
Thus, for one integer $c$,
\[
 2a_1=c,\qquad
 a_{i-1}+2a_i=c\quad(2\le i<n),\qquad
 a_{n-1}+3a_n=c.
 \tag{P-F19k}
\]
Adding these equations gives $nc=3D$, so $c=3s/g$.
The first $n-1$ equations solve to
\[
 a_i=\frac{c\bigl(2^i-(-1)^i\bigr)}{3\,2^i}
       \qquad(1\le i<n),
 \tag{P-F19l}
\]
as is checked by substitution and the first equation.
The numerator factor $2^i-(-1)^i$ is odd. Since $a_{n-1}$
is an integer, (P-F19l) forces $2^{n-1}\mid c$.
But $c=(3/g)s$ and $3/g$ is odd, so $2^{n-1}\mid s$.
This proves (P-F19j). Possible cancellation among the surviving
monomials causes no problem: nonzero evaluation guarantees that
at least one such monomial occurs, and every one obeys the same
divisibility condition.

\emph{Diagonal-stabilizer calculation.}
Let $\zeta$ be a primitive root of unity of order $N=3\,2^{n-1}$
and define
\[
 A=\operatorname{diag}
  (\zeta,\zeta^{-2},\zeta^{(-2)^2},\ldots,
                         \zeta^{(-2)^{n-1}})\in GL_n.
\]
Each chain monomial in (P-F19h) is fixed by $A$, since
$2(-2)^{i-1}+(-2)^i=0$; the last monomial is fixed since
$3(-2)^{n-1}$ is a multiple of $N$. The same is true with
$A^{-1}$, so $A\cdot F_n=F_n$ for the specified action on forms.
Set
\[
 S_n=1-2+(-2)^2+\cdots+(-2)^{n-1}
     =\frac{1-(-2)^n}{3}.
\]
The integer $S_n$ is odd and satisfies $S_n\equiv n\pmod3$.
Consequently
\[
 \operatorname{ord}(\det A)
 =\frac{N}{\gcd(N,S_n)}
 =\frac{3\,2^{n-1}}{g}.
 \tag{P-F19m}
\]
An $SL_n$-invariant polynomial $P$ of coefficient degree
$D=d_0s$ satisfies, as a function on $W^*$,
\[
 P(a\cdot F)=(\det a)^{-3s/g}P(F).
 \tag{P-F19n}
\]
This follows by writing $a=zh$ as above: $h$ does not change
$P$, while a scalar sends $F$ to $z^{-3}F$.
If $P(F_n)\ne0$, evaluation at its stabilizer $A$ therefore
requires $(\det A)^{3s/g}=1$. By (P-F19m),
$(3\,2^{n-1}/g)\mid(3s/g)$, which is again equivalent to
$2^{n-1}\mid s$. This verifies (P-F19j) independently of
the coefficient recurrence.

Finally, let $q_1,\ldots,q_r$ be any homogeneous system of
parameters of $T^{(n)}$. The algebra is finite as a module over
$\mathbb C[q_1,\ldots,q_r]$, so
$T^{(n)}/(q_1,\ldots,q_r)$ is finite-dimensional and graded.
Every positive-degree homogeneous element in this quotient is
nilpotent: its powers eventually have degree beyond the largest
nonzero graded piece. In particular
$\Delta^M\in(q_1,\ldots,q_r)$ for some integer $M\ge1$.
If every $q_i(F_n)$ were zero, evaluation of this ideal membership
would give $\Delta(F_n)^M=0$, contradicting (P-F19i).
Thus some $q_i$ is nonzero at $F_n$. Its positive stretch degree
is divisible by $2^{n-1}$ by (P-F19j), and is at least that
large. This applies to every h.s.o.p., including one with a
lexicographically minimal degree sequence, proving (P-F19c).

The inner module is the single irreducible $S^{(3)}V$, and
the partitions in (P-F19a) have heights one and $n$.
Thus (P-F19d) follows directly from the source's stated
parameter. Exponential growth of $2^{n-1}$ exceeds every
polynomial in $h=O(n^6)$, proving the asserted failure of the
canonical modular-index clause.
\end{proof}

\paragraph{Exact scope of the refutation.}
GCT VI, Hypothesis 1.6.7 first asks for some positive rational
form of polynomial modular index, and then explicitly asks for
this bound on the minimal form arising from a lexicographically
minimal h.s.o.p. of its canonical ring. P-F19 refutes the latter,
stronger clause. Its proof concerns values of actual invariants
at a particular semistable point. A positive rational expression
for the Hilbert series need not arise from parameters of that
ring, so the first existence clause is not refuted by this
argument alone. P-F30--P-F31 subsequently close that gate by proving
an actual primitive-$2^{n-1}$ pole of the ordinary Hilbert series,
and thereby refute the first clause as well. The paired-wave argument below supplies the additional
estimates needed to prove superpolynomial positivity and saturation indices.

This is an irreducible-inner ordinary plethysm family, unlike
the reducible torus embeddings in P-F14--P-F15. Determinant
twisting is absent: the inner shape is always $(3)$, and only
the number of variables grows. No historical
priority is asserted for the chain forms, their stabilizers,
or the resulting invariant-degree obstruction. The same chain
forms, after reversing the variables, already appear in
Bus\'e--Jouanolou, Example 4.12, printed p.~42, with attribution
to Demazure; the statement here gives their explicit consequence
for the specified canonical plethysm parameter bound.
